% Generated from locale/tr/content/methods/induction/relations.tex; do not hand-edit.


\olfileid[tr]{mth}{ind}{rel}

\olsection{Bağıntılar ve Fonksiyonlar}

Bir nesneler kümesini (doğal sayılar ya da düzgün terimler gibi)
tümevarımsal olarak tanımladıktan sonra, bu nesneler \emph{üzerindeki
bağıntıları} da tümevarımla tanımlayabiliriz. Örneğin şu düşünceyi ele
alın: bir düzgün~$t_1$ terimi, bir düzgün~$t_2$ teriminin parçası olarak
geçiyorsa onun altterimidir. Bunun için $t_1 \sqsubseteq t_2$ simgesini
kullanalım. Elbette her düzgün terim kendisinin bir altterimidir: $t
\sqsubseteq t$. Bu bağıntının tümevarımsal tanımını şöyle verebiliriz:

\begin{defn}
  Bir düzgün~$t_1$ teriminin~$t_2$'nin altterimi olması bağıntısı, $t_1
  \sqsubseteq t_2$,~$t_2$ üzerinde tümevarımla şöyle tanımlanır:
  \begin{enumerate}
  \item $t_2$ bir harfse $t_1 \sqsubseteq t_2$, ancak ve ancak $t_1 = t_2$
    ise geçerlidir.
  \item $t_2$, $[s_1 \circ s_2]$ ise $t_1 \sqsubseteq t_2$, ancak ve
    ancak $t_1 = t_2$, $t_1 \sqsubseteq s_1$ ya da $t_1 \sqsubseteq s_2$
    ise geçerlidir.
  \end{enumerate}
\end{defn}

Bu tanım bize örneğin $\mathrm{a} \sqsubseteq [\mathrm{b} \circ
\mathrm{a}]$ olduğunu söyler. Çünkü (2), $\mathrm{a} \sqsubseteq
[\mathrm{b} \circ \mathrm{a}]$'nın ancak ve ancak $\mathrm{a} =
[\mathrm{b} \circ \mathrm{a}]$ ya da $\mathrm{a} \sqsubseteq
\mathrm{b}$ ya da $\mathrm{a} \sqsubseteq \mathrm{a}$ ise geçerli
olduğunu söyler. İlk ikisi yanlıştır: $\mathrm{a}$ açıkça $[\mathrm{b}
\circ \mathrm{a}]$ ile özdeş değildir ve (1)'e göre $\mathrm{a}
\sqsubseteq \mathrm{b}$ ancak ve ancak $\mathrm{a} = \mathrm{b}$ ise
geçerlidir; bu da yanlıştır. Bununla birlikte yine (1)'e göre,
$\mathrm{a} \sqsubseteq \mathrm{a}$ ancak ve ancak $\mathrm{a} =
\mathrm{a}$ ise geçerlidir; bu doğrudur.

Bu tanımın başarısının henüz kanıtlamadığımız bir olguya dayandığına
dikkat etmek önemlidir: her düzgün~$t$ terimi ya tek başına bir harftir
ya da \emph{biricik biçimde belirlenmiş} düzgün $s_1$ ve $s_2$ terimleri
için $t = [s_1 \circ s_2]$ olur. Buradaki ``biricik biçimde belirlenmiş'',
$t = [s_1 \circ s_2]$ iken \emph{aynı zamanda} $= [r_1 \circ r_2]$ olup
$s_1 \neq r_1$ ya da $s_2 \neq r_2$ olamayacağı anlamına gelir. Durum
böyle olsaydı (2). madde kendi kendisiyle
çatışabilirdi: $t_2$'yi $[s_1 \circ s_2]$ olarak okuyunca $t_1
\sqsubseteq t_2$ sonucunu, fakat $t_2$'yi $[r_1 \circ r_2]$ olarak
okuyunca $t_1 \sqsubseteq t_2$ olmadığını elde edebilirdik. Bunun
olamayacağını kanıtlamadan önce, \emph{olabildiği} bir örneğe bakalım.

\begin{defn}
  \emph{Parantezsiz terimleri} tümevarımsal olarak şöyle tanımlayın:
  \begin{enumerate}
  \item Her harf bir parantezsiz terimdir.
    \item $s_1$ ile $s_2$ parantezsiz terimlerse $s_1 \circ s_2$ de
      parantezsiz bir terimdir.
    \item Başka hiçbir şey parantezsiz terim değildir.
  \end{enumerate}
\end{defn}

Parantezsiz terimlere örnek olarak $\mathrm{a}$, $\mathrm{b} \circ
\mathrm{d}$, $\mathrm{b} \circ \mathrm{a} \circ \mathrm{b}$ verilebilir.
Parantezsiz terimler için ``altterim''i yukarıdaki gibi tanımlasaydık,
ikinci madde şöyle olurdu:
\begin{center}
  $t_2 = s_1 \circ s_2$ ise $t_1 \sqsubseteq t_2$, ancak ve ancak $t_1
  = t_2$, $t_1 \sqsubseteq s_1$ ya da $t_1 \sqsubseteq s_2$ ise
  geçerlidir.
\end{center}

Şimdi $\mathrm{b} \circ \mathrm{a} \circ \mathrm{b}$, şu değerlerle
$s_1 \circ s_2$ biçimindedir:
\begin{align*}
  s_1 & = \mathrm{b} \text{ ve} & s_2 & = \mathrm{a} \circ \mathrm{b}.
\intertext{Aynı zamanda şu değerlerle $r_1 \circ r_2$ biçimindedir:}
  r_1 & = \mathrm{b} \circ \mathrm{a} \text{ ve} & r_2 &= \mathrm{b}.
\end{align*}
Peki $\mathrm{a} \circ \mathrm{b}$, $\mathrm{b} \circ \mathrm{a}
\circ \mathrm{b}$'nin bir altterimi midir? İlk okuyuşa göre yanıt evet,
ikinci okuyuşa göre hayırdır.

Bir düzgün terimin başka düzgün terimlerden kurulma yolunun tek türlü olması
özelliğine \emph{tek türlü okunabilirlik} denir. Böyle tümevarımsal olarak
tanımlanan nesneler için bağıntıların tümevarımsal tanımları önemli
olduğundan, bu özelliğin geçerli olduğunu kanıtlamalıyız.

\begin{prop}
  $t$'nin bir düzgün terim olduğunu varsayın. O zaman ya $t$ tek başına
  bir harftir ya da biricik biçimde belirlenmiş düzgün $s_1$, $s_2$
  terimleri için~$t = [s_1 \circ s_2]$ olur.
\end{prop}

\begin{proof}
  $t$ tek başına bir harfse koşul sağlanır. Öyleyse $t$'nin tek başına
  bir harf olmadığını varsayın. Tümevarımsal tanımdan, bu durumda
  $t$'nin bazı düzgün $s_1$ ve~$s_2$ terimleri için $[s_1 \circ s_2]$
  biçiminde olması gerektiğini görebiliriz. Geriye bunların biricik biçimde
  belirlendiğini, yani $t = [r_1 \circ r_2]$ ise $s_1 = r_1$ ve $s_2 =
  r_2$ olduğunu göstermek kalır.

  Öyleyse düzgün $s_1$, $s_2$, $r_1$, $r_2$ terimleri için $t = [s_1
  \circ s_2]$ ve aynı zamanda $t = [r_1 \circ r_2]$ olduğunu varsayın.
  $s_1 = r_1$ ve $s_2 = r_2$ olduğunu göstermeliyiz. Önce $s_1$ ile
  $r_1$ özdeş olmalıdır; yoksa biri ötekinin öz başlangıç parçası olur.
  Fakat \olref[sti]{prop:initial} gereği $s_1$ ile~$r_1$'in ikisi de
  düzgün terimse bu olanaksızdır. $s_1 = r_1$ ise açıkça $s_2 = r_2$ de
  geçerlidir.
\end{proof}

Fonksiyonları da tümevarımsal olarak tanımlayabiliriz: örneğin herhangi
bir düzgün terimi, içindeki iç içe $[\dots]$ yapısının en büyük
derinliğine gönderen~$f$ fonksiyonunu şöyle tanımlayabiliriz:

\begin{defn}
  \ollabel{defn:depth} Bir düzgün terimin \emph{derinliği}, $f(t)$,
  tümevarımsal olarak şöyle tanımlanır:
  \[
  f(t) = \begin{cases}
    0 & \text{ $t$ bir harfse}\\
    \max(f(s_1), f(s_2)) + 1 & \text{ $t = [s_1 \circ s_2]$ ise.}
  \end{cases}
  \]
\end{defn}

Örneğin
\begin{align*}
  f([\mathrm{a} \circ \mathrm{b}]) & =
    \max(f(\mathrm{a}),f(\mathrm{b})) + 1 = \\
  &= \max(0, 0) + 1 = 1, \text{ ve}\\
  f([[\mathrm{a} \circ \mathrm{b}] \circ \mathrm{c}]) & =
    \max(f([\mathrm{a} \circ \mathrm{b}]), f(\mathrm{c})) + 1 = \\
  & = \max(1,0) + 1 = 2.
\end{align*}

Burada elbette $s_1$ ile $s_2$'nin düzgün terimler olduğunu varsayar ve
her düzgün terimin ya bir harf ya da $[s_1 \circ s_2]$ biçiminde olduğu
olgusunu kullanırız. Bu biçimde yalnızca tek yolla olabilmesi yine
önemlidir. Nedenini görmek için daha önce tanımladığımız parantezsiz
terimleri yeniden ele alın. Buna karşılık gelen ``tanım'' şöyle olurdu:
\[
  g(t) =
  \begin{cases}
    0 & \text{ $t$ bir harfse}\\
   \max(g(s_1), g(s_2)) + 1 & \text{ $t = s_1 \circ s_2$ ise.}
  \end{cases}
\]
Şimdi $\mathrm{a} \circ \mathrm{b} \circ \mathrm{c} \circ
\mathrm{d}$ parantezsiz terimini ele alın. Birden fazla biçimde
okunabilir; örneğin şu değerlerle $s_1 \circ s_2$ olarak:
\begin{align*}
  s_1 & = \mathrm{a} \text{ ve} &
  s_2 & = \mathrm{b} \circ \mathrm{c} \circ \mathrm{d},
\intertext{ya da şu değerlerle $r_1 \circ r_2$ olarak:}
  r_1 & = \mathrm{a} \circ \mathrm{b} \text{ ve} &
  r_2 &= \mathrm{c} \circ \mathrm{d}.
\end{align*}
$g$'yi ilk okuyuşa göre hesaplamak şunu verirdi:
\begin{align*}
  g(s_1 \circ s_2) & =
  \max(g(\mathrm{a}), g(\mathrm{b} \circ \mathrm{c} \circ \mathrm{d})) + 1 =\\ & =
  \max(0,2) + 1 = 3
  \intertext{öteki okuyuşa göre ise şunu elde ederiz:}
  g(r_1 \circ r_2) & =
  \max(g(\mathrm{a} \circ \mathrm{b}), g(\mathrm{c} \circ \mathrm{d})) + 1=\\ &
  = \max(1,1) + 1 = 2
\end{align*}
Fakat bir fonksiyon her zaman tek bir değer vermelidir; öyleyse~$g$
``tanımımız'' aslında hiçbir fonksiyon tanımlamaz.

\begin{prob}
  $l(t)$, düzgün~$t$ terimindeki simgelerin sayısı olmak üzere, $l$
  fonksiyonunun tümevarımsal bir tanımını verin.
\end{prob}

\begin{prob}
  Düzgün terimler~$t$ üzerinde yapısal tümevarımla $f(t) < l(t)$
  olduğunu kanıtlayın (burada $l(t)$,~$t$'deki simgelerin sayısı ve
  $f(t)$, \olref[mth][ind][rel]{defn:depth}'te tanımlandığı üzere
  $t$'nin derinliğidir).
\end{prob}

